% Part: incompleteness
% Chapter: representability-in-q
% Section: prim-rec

\documentclass[../../../include/open-logic-section]{subfiles}

\begin{document}

\olfileid{inc}{req}{pri}
\olsection{આદિમ પુનરાવર્તનનું અનુસરણ}

હવે બીટા વિધેય વડે આદિમ
પુનરાવર્તનથી થતી વ્યાખ્યાનું
નિયમિત લઘુત્તમીકરણથી ``અનુસરણ''
કરી શકાય છે તે બતાવી શકીએ.
માનો કે આપણી પાસે $f(\vec x)$
અને $g(\vec x, y, z)$ છે.
તો $f$ અને~$g$માંથી આદિમ
પુનરાવર્તન વડે વ્યાખ્યાયિત વિધેય
$h(\vec x,y)$ આ છે:
\sourcecorrection{OLINC-025}{સ્થિર અંગ્રેજી
મૂળના આ ગદ્યમાં $h(x,\vec z)$
લખાયું છે, જ્યારે તરત પછીનાં
બંને સમીકરણો અને બાકીની
સાબિતીમાં $h(\vec x,y)$ છે.
અહીં ચલોના એ જ ક્રમનો ઉપયોગ
કર્યો છે.}
\begin{align*}
h(\vec x, 0) & =  f(\vec x) \\
h(\vec x, y+1) & =  g(\vec x, y, h(\vec x, y)).
\end{align*}
$f$ અને~$g$માંથી માત્ર સંયોજન
અને નિયમિત લઘુત્તમીકરણ વડે~$h$
વ્યાખ્યાયિત થાય છે તે બતાવવું છે.
એ માટે મૂળભૂત વિધેયો અને
તેમનામાંથી સંયોજન તથા નિયમિત
લઘુત્તમીકરણ વડે વ્યાખ્યાયિત
વિધેયો (જેમ કે~$\beta$) વાપરીશું.

\begin{lem}
\ollabel{lem:prim-rec}
$h$ને $f$ અને $g$માંથી આદિમ
પુનરાવર્તન વડે વ્યાખ્યાયિત
કરી શકાય, તો $f$, $g$ અને
વિધેયો $\Zero$, $\Succ$,
$\Proj{n}{i}$, $\Add$, $\Mult$,
$\Char{=}$માંથી સંયોજન અને
નિયમિત લઘુત્તમીકરણ વડે પણ
વ્યાખ્યાયિત કરી શકાય છે.
\end{lem}

\begin{proof}
પહેલાં સહાયક વિધેય $\hat h(\vec x, y)$
વ્યાખ્યાયિત કરો. તે નીચેની
શરતો સંતોષતી શ્રેણીને સંકેતિત
કરતી ન્યૂનતમ સંખ્યા~$d$ આપે છે;
આ~$d$ શ્રેણીનો સંકેત છે:
\begin{enumerate}
\item $(d)_0 = f(\vec x)$, અને
\item દરેક $i < y$ માટે $(d)_{i+1} = g(\vec x, i, (d)_i)$,
\end{enumerate}
અહીં હવે $(d)_i$ એ $\beta(d,i)$નું
ટૂંકું લખાણ છે. બીજા શબ્દોમાં,
$\hat h$ એ
$\tuple{h(\vec x, 0), h(\vec x, 1),
\dots, h(\vec x, y)}$ શ્રેણીનો
સંકેત આપે છે.
\sourcecorrection{OLINC-026}{સ્થિર અંગ્રેજી
મૂળમાં $\hat h$ શ્રેણી પોતે
આપે છે એવું કહેવાયું છે;
પરંતુ અગાઉની વ્યાખ્યા મુજબ
તે શ્રેણીનો સંકેત~$d$ આપે છે.
અહીં એ ભેદ સ્પષ્ટ કર્યો છે.}
$\hat h$ને આ પ્રમાણે લખી શકીએ:
\[
\hat h(\vec x, y) = \umin{d}{(\beta(d,0) = f(\vec x) \land \bforall{i <
  y}{\beta(d,i+1) = g(\vec x, i,\beta(d,i)})}.
\]
ધ્યાન આપો કે અહીં આદિમ
પુનરાવર્તન જરૂરી નથી, માત્ર
લઘુત્તમીકરણ પૂરતું છે. બીટા
વિધેયના લેમ્મા
\olref[bet]{lem:beta}ને કારણે
લઘુત્તમીકરણ હેઠળનું વિધેય
નિયમિત છે.

હવે, આ મળે છે:
\[
h(\vec x, y) = \beta(\hat h(\vec x, y), y),
\]
આથી~$h$ને મૂળભૂત વિધેયોમાંથી
માત્ર સંયોજન અને નિયમિત
લઘુત્તમીકરણ વડે વ્યાખ્યાયિત
કરી શકાય છે.
\end{proof}

\end{document}
